\section{Methodology}
\subsection{Formulation}
The system consists of $N$ omnidirectional mobile robots maninupulating an object with polygonal shape whose motion is constrained to a 2D plane using pushing interaction. The primary objective is for the robots to cooperatively perform the estimation task to obtain the Center of Mass (COM) of the object in a distributed manner. It is assumed that the object's position and 2D geometry is known to all robot and neither the object nor the robots are subject to kinematic constraints. Each robot $i$ has a circular footprint with radius $r_R$, can exert a normal force $f_i$ within a limited range : $0 < f_i < f_{max}$, as well as measure the position $p_i$ and velocity $\dot{p}_i$ at its contact point $C_i$ on the object's surface. The forces applied to the object by the robots are represented by the vector: $\mathbf{f}=[f_1,\ldots,f_N]^T \in
\mathcal{F}
:=
\left\{
\mathbf{f}\in\mathbb{R}^N
\mid
0\leq f_i\leq f_{\max},
\; i=1,\ldots,N
\right\}$. It is also assumed that the robots are able to maintain continuous contact with the object and control an accurately regulated pushing force at its corresponding contact point throughout the manipulation process. The net force and torque that all robots exert on the object in the world-fix frame $W$ are descibe as:

\begin{equation}
F
=
\sum_{i=1}^{N}
\mathbf{F}_{i}
\label{eq:force}
\end{equation}

\begin{equation} 
\tau 
= 
\sum_{i=1}^{N} 
\left( 
\mathbf{p}_{i}-\mathbf{p}_{C} 
\right)^{\perp T} 
\mathbf{F}_{i} 
\label{eq:tau}
\end{equation}

where $(\cdot)^\perp$ denotes the perpendicular operator, defined as
$\mathbf{v}^\perp = Q\mathbf{v}$, with
$Q=\begin{bmatrix}0&-1\\1&0\end{bmatrix}$. $F_i = f_i ^{\perp T}n_i$ is the force vector of robot $i$, with $n_i$ being the unit normal vector at the contact point $C_i$. The friction model is negligible since we do not require a dynamically accurate model. The equation \eqref{eq:tau} can be rewritten as:

\begin{equation} 
\tau 
= 
\sum_{i=1}^{N} 
\mathbf{p}_{i}^{\perp T}\mathbf{F}_{i} 
- 
\mathbf{p}_{C}^{\perp T} \mathbf{F}
\label{eq:tau2}
\end{equation} 

\subsection{Estimation algorithm}

When the object undergoes rotational motion, all contact points rotate instantaneously about the same instantaneous center of rotation (ICR). The velocity of each contact point is tangential to its instantaneous circular trajectory and therefore perpendicular to the line connecting the contact point to the ICR. Consequently, the ICR can be geometrically determined from the velocity measurements at multiple contact points by the equation below:
\begin{equation}
\mathbf{p}_{ICR}
=
\mathbf{p}_{i}
-
\mathbf{p}_{o}
+
\frac{1}{\omega_o}
\left(\dot{\mathbf{p}}_{i}\right)^{\perp}
\end{equation}

where $p_o$ and $\omega_o$ are the position and angular velocity of the object, respectively.

The proposed distributed estimation algorithm idea is to generate a pure rotational motion of the object by controlling formation and the exert force of the robots. When this occurs, the ICR is located exactly at the COM of the object. Note that since the net force $F$ is independent of the object's COM position, when the object undergoes pure rotation, in other words, the net force $F$ on the object is zero, then \eqref{eq:tau2} reduces to:

\begin{equation} 
\tau 
= 
\sum_{i=1}^{N} 
\mathbf{p}_{i}^{\perp T}\mathbf{F}_{i}
\end{equation}

Therefore, the net torque $\tau$ is only dependent on the contact points position and the exerted forces, regardless of the object's COM position. By controlling the formation of the robots and their exerted forces, we can achieve the desired net torque $\tau_d$ that will induce a pure rotational motion of the object. This torque can be agreed upon by all robots through communication or a consensus algorithm. The COM position can then be estimated by each robot using the following equation:

\begin{equation}
{}^{B}\mathbf{p}_{\mathrm{C}}
=
R^{T}(\theta)
\left[
{}^{W}\mathbf{p}_{i}
-
{}^{W}\mathbf{p}_{o}
+
\frac{1}{\omega}
\left({}^{W}\dot{\mathbf{p}}_{i}\right)^{\perp}
\right]
\end{equation}

The estimation problem can now be formulated into 2 sub-problems: formation synthesis and force-based control. The formation synthesis problem's objective is to find a configuration $p$,  which is able to cooperatively exert the desired force and torque on the object. The force control problem determine the forces to be exerted by each robot to achieve the desired net force and torque on the object. 

\subsection{Formation synthesis}

To efficiently search for the optimal robot configuration, the search space is defined directly along the object boundary. Since each robot is constrained to establish contact with the object's surface, the position of robot i can be parameterized by a scalar variable $s_i$ representing the arc length along the polygonal boundary of the object. Accordingly, the search space for (N) robots is expressed as $\mathbf{s}=[s_1,\ldots,s_N]^T$, where $s_i\in[0,L)$ and (L) denotes the total perimeter of the object. To improve contact stability, a small margin $d_s$ around each polygon vertex is omitted from the search space, the resulting feasible search space is illustrated in.

Since the maximum pushing force $f_{\max}$ is finite due to the physical limitations of the robots, for each configuration of the robots, the corresponding manipulation space is a closed set $W$ consisting of all feasible set of forces and torques, or wrenches that can be exerted by the current formation. The full manipulation space can not be accurately determined or analyzed without prior information about the COM position, which heavily affects the rotational motion of the object. However as established above, when pure rotation occurs, the torque becomes decoupled from the COM position. The maximum and minimum torque can be exerted by a configuration $p$ can be derived by solving the following linear programming problems:

\begin{equation}
\begin{aligned}
\tau_{\max}(\mathbf{p})
&=
\max_{\mathbf{f}\in\mathcal{F}}
\sum_{i=1}^{N}
\mathbf{p}_i^{\perp T}\mathbf{F}_i
\quad
\text{s.t.}\quad
\sum_{i=1}^{N}\mathbf{F}_i=\mathbf{0},
\\
\tau_{\min}(\mathbf{p})
&=
\min_{\mathbf{f}\in\mathcal{F}}
\sum_{i=1}^{N}
\mathbf{p}_i^{\perp T}\mathbf{F}_i
\quad
\text{s.t.}\quad
\sum_{i=1}^{N}\mathbf{F}_i=\mathbf{0}.
\end{aligned}
\end{equation}

These two quantities are, in fact, the intersection of the manipulation space with the torque axis. A formation is considered feasible if the desired torque $\tau_d$ is within the range of $[\tau_{\min}, \tau_{\max}]$. It is also benificial to maximize the minimum distance between the desired torque $\tau_d$ and the boundary of the feasible torque range, which can be formulated as a fitness function $J_r$ for robustness metric:

\begin{equation}
J_r
=
\min
\left(
\left|\tau_{\max}-\tau_d\right|,
\left|\tau_{\min}-\tau_d\right|
\right).
\end{equation}

The formation synthesis problem is formulated as:

\begin{subequations}
\begin{align}
\min_{\mathbf{p} \in \mathbb{R}_{\geq 0}^{N}}
\quad & J(\mathbf{p}) = -J_r(\mathbf{p})
\label{eq:objective} \\[2mm]
\text{s.t.}\quad
& \exists\,\tilde{\mathbf{f}} \in \mathcal{F}
\quad \text{s.t.}\quad
F = \mathbf{0},\;
\sum_{i=1}^{N} \mathbf{p}_{i}^{\perp T}\mathbf{F}_{i}
= \tau_{d},
\label{eq:wrench_constraint} \\[2mm]
& \mathrm{frD}(p_i,p_j) \geq d_{\min},
\qquad \forall i \neq j .
\label{eq:distance_constraint}
\end{align}
\label{eq:formation_synthesis}
\end{subequations}

where $d_{min} > 2r_r$ is the a minimum Euclidean distance between the robots to avoid collision. It is essential that the constraints \eqref{eq:wrench_constraint} and \eqref{eq:distance_constraint} are satisfied to ensure that the desired wrench can be exerted by the formation without violating physical feasibility, while minimizing the objective function yields a configuration that is optimal with respect to the selected criteria. The objective function \eqref{eq:objective} is a desirable property, which can be modified to include other metrics. In this case, only rotational robustness is consisered. The formation synthesis problem is a non-convex optimization problem, which can be solved using heuristic approaches, for example Particle Swarm Optimization (PSO). To keep the distributed nature of the system, each robot can run the optimization algorithm independently and share their best particle with their neighbors. The best solution is then selected as the final configuration.

\subsection{Force control}

The next problem to be solved is to determine the forces to be exerted by each robot to achieve the desired net force and torque on the object. While one naive approach would be to just solve \eqref{eq:force} and \eqref{eq:tau2} using least-square approaches, this would fail to account for the physical limitations of the forces exerted by the robots, resulting in a solution that is not physically feasible. To tackle this issue, the force control problem is formulated as an optimization problem:

\begin{align}
\min_{\mathbf{f}} \quad
& \left(
\sum_{i=1}^{n}
\mathbf{p}_{i}^{\perp T}\mathbf{F}_{i}
-\tau_d
\right)^2
\label{eq:objective}\\
\text{s.t.}\quad
& \mathbf{F}=\mathbf{0},
\label{eq:force_balance}\\
& 0 \leq f_i\leq f_{\max},
\qquad i=1,\ldots,n.
\label{eq:force_bound}
\end{align}

Satisfying constraints \eqref{eq:force_balance} and \eqref{eq:force_bound} is the primary requirement of the optimization problem, since they ensure the physical feasibility and enable the desired pure rotational motion. It is desirable that the net torque is as close as possible to the desired torque $\tau_d$, howerver it is not required to exactly match $\tau_d$, provided that its magnitude is sufficient to overcome the resisting static friction torque.